% Slide type: plain
\begin{frame}{Section 5.2A ─ バンドコンビネーション}
\textbf{CA の 3 形態}
\begin{itemize}
  \item Intra-band Contiguous: 同一バンド内の隣接 CC（RF 実装が最も単純）
  \item Intra-band Non-contiguous: 同一バンド内、間に他社帯域あり
  \item Inter-band: 異なるバンドの CC（実運用で最も一般的）
\end{itemize}

\textbf{表記法}: \texttt{CA\_n1A-n78C} = n1 で 1CC + n78 で 2CC（クラス C, 100--200 MHz）

\textbf{5.2 系列の拡張}
\begin{itemize}
  \item 5.2A=CA, 5.2B=DC, 5.2C=SUL
  \item 5.2D=UL MIMO, 5.2K=多方向 Rx/Tx, 5.2M=LP-WUS
\end{itemize}
\end{frame}
